% Generate Less, Classify More
% ICLR 2027 submission draft. Double-blind: NO author names/affiliations.
% NOTE: swap this preamble for the official iclr2027_conference.sty when it is
% released; the article-based fallback below lets the draft compile now.
\documentclass{article}
\usepackage[utf8]{inputenc}
\usepackage{times}
\usepackage{graphicx}
\usepackage{booktabs}
\usepackage{amsmath,amssymb}
\usepackage{hyperref}
\usepackage[margin=1in]{geometry}
\usepackage{multirow}
\usepackage{natbib}

\title{Generate Less, Classify More: Probe-Guided Pruning and\\
Discriminative Heads for Efficient Intent Detection and Slot Filling\\
with Small Language Models}
\author{Anonymous authors\\Paper under double-blind review}
\date{}

\begin{document}
\maketitle

\begin{abstract}
Task-oriented dialogue increasingly relies on small language models (SLMs) that
emit structured output (intent + slots) via autoregressive generation. We show
that for \emph{closed-set} intent detection and slot filling this is both slow
and brittle, and unnecessary. We present a recipe that converts a generative SLM
into a \textbf{discriminative single-pass} model: (1) a frozen linear-probe
layer sweep selects a pruned backbone depth \emph{a priori}, before any
fine-tuning; (2) a linear intent head and a token-level slot head (softmax or
CRF) replace token-by-token generation. On a 124M-parameter decoder, pruning to
3 of 12 layers, our model runs at $\sim$5\,ms P50 on a commodity CPU versus
$\sim$359\,ms for the generative baseline (a $\sim$70$\times$ speedup on ATIS),
improves intent accuracy from 48.3\% to 97.95\%, and eliminates a generative
parse-failure rate of up to 46.7\% (0\% by construction). A 15-second probe
correctly selects a 4$\times$-shallower depth on all five of ATIS, SNIPS,
MASSIVE, CLINC150, and BANKING77; across three seeds the selected depth is
statistically indistinguishable from full depth on three datasets and within
$\sim$1--2.6 points on the two most fine-grained. We further contribute a
negative-results study of CPU inference optimizations: dynamic INT8 and
\texttt{torch.compile} give no benefit for shallow models, while ONNX Runtime
yields a $4.5\times$ speedup at the cost of brittle export.
\end{abstract}

\section{Introduction}
Small language models are increasingly deployed for task-oriented natural
language understanding (NLU), where a user utterance is mapped to a discrete
\emph{intent} and a set of typed \emph{slots}. A common pattern is to fine-tune
a generative SLM to emit a structured object (e.g.\ JSON
\texttt{\{"intent":...,"slots":...\}}) autoregressively. This inherits two costs
from generation: (i) latency that scales with output length, and (ii) a nonzero
rate of \emph{malformed} outputs that fail to parse into the target schema.

We argue that closed-set intent + slot filling does not need generation. Our
central claim: \emph{converting a generative SLM into a probe-pruned
discriminative single-pass model yields a large CPU speedup at equal-or-better
accuracy and zero structured-output parse failures, and the pruned depth can be
chosen a priori from a cheap linear probe.}

\textbf{Contributions.}
\textbf{(C1)} \emph{Probe-guided depth selection}: a linear probe on frozen
hidden states predicts a near-optimal pruned depth before expensive fine-tuning.
\textbf{(C2)} \emph{Generative$\rightarrow$discriminative conversion}: a large
speedup plus a reliability guarantee (parse-failure rate $\rightarrow$ 0 by
construction). \textbf{(C4)} a negative-results study of common CPU inference
optimizations and their root causes. We position our contribution honestly:
against the generative baseline the win is decisive; against a purpose-built
bidirectional encoder we are competitive, and our recipe requires no separate
encoder---it reuses an existing decoder.

\section{Related Work}
\textbf{Joint intent/slot models.} JointBERT~\citep{chen2019bert} and successors
fine-tune an encoder with intent and slot heads. We revisit this from the angle
of \emph{converting a decoder} and \emph{choosing depth by probing}.
\textbf{Structured/constrained decoding}~\citep{willard2023guidance} makes
generation valid but retains autoregressive cost; discriminative heads make
validity structural at no decode cost.
\textbf{Compression}: layer dropping~\citep{fan2020layerdrop}, early
exit~\citep{xin2020deebert}, and rotation-based quantization~\citep{ashkboos2024quarot}.
We use static depth pruning guided by probes rather than dynamic exit.
\textbf{Probing}~\citep{alain2017probes,tenney2019bert} analyzes what layers
encode; we repurpose probes as an \emph{a-priori depth-selection procedure}.

\section{Method}
\textbf{Setup.} A closed intent set $\mathcal{I}$ and BIO slot tags
$\mathcal{T}$; we predict one intent and a tag per token in a single forward
pass. \textbf{(3.1) Probe-guided depth selection.} We freeze the pretrained
decoder and, in one forward pass with hidden-state output, mean-pool the hidden
states at candidate depths $\{L/4, L/2, 3L/4, L\}$ over non-pad tokens. A linear
probe (logistic regression) is trained on each depth's pooled features to
predict the intent. We select the shallowest depth whose probe accuracy is
within $\epsilon$ of the best (we use $\epsilon=0.02$). \textbf{(3.2)
Discriminative heads.} We truncate the backbone to the selected depth, add a
linear intent head over mean-pooled hidden states and a per-token slot head
(softmax cross-entropy, or a linear-chain CRF~\citep{lafferty2001crf} with
Viterbi decoding). The joint loss is
$\mathcal{L} = \mathcal{L}_{\text{intent}} + \lambda\,\mathcal{L}_{\text{slot}}$
with $\lambda=2$; we length-normalize the CRF NLL so it does not dominate the
intent term. \textbf{(3.3) Reliability.} Because output is a fixed set of
classifier decisions, it is structurally valid: parse-failure rate is 0 by
construction.

\section{Experimental Setup}
\textbf{Datasets.} ATIS~\citep{hemphill1990atis},
SNIPS~\citep{coucke2018snips}, MASSIVE-en~\citep{fitzgerald2023massive}
(intent+slot); CLINC150~\citep{larson2019clinc} and
BANKING77~\citep{casanueva2020banking77} (intent-only, 150 and 77 classes).
\textbf{Models.} Backbone: gpt2-124M~\citep{radford2019gpt2}; encoder baseline:
DistilBERT~\citep{sanh2019distilbert}. \textbf{Metrics.} Intent accuracy,
span-level slot F1, parse-failure rate, and single-example P50 latency on a
commodity CPU (Apple Silicon arm64, 10 threads). For fairness across the
architecture comparison we match training-set size (5000 examples) between the
encoder baseline and the pruned model. The generative baseline is bounded to
3000 train / 300 eval for CPU-tractable autoregressive decoding.

\section{Results}

\begin{table}[t]\centering
\caption{Main comparison on ATIS. The pruned discriminative model is
$\sim$70$\times$ faster than the generative baseline, matches encoder intent
accuracy, and never emits a malformed output.}
\begin{tabular}{lrrrr}
\toprule
Model & Intent Acc & Slot F1 & Parse-fail & P50 (ms)\\
\midrule
Generative SLM (gpt2, autoregressive) & 48.3 & 44.6 & 46.7 & 359\\
Encoder discriminative (DistilBERT) & \textbf{98.98} & \textbf{95.1} & 0 & 7.0\\
\textbf{Ours: probe-pruned (gpt2 depth-3)} & 97.95 & 91.0 & \textbf{0} & \textbf{5.1}\\
\bottomrule
\end{tabular}
\end{table}

\textbf{5.1 Probe-guided depth (C1).} On ATIS the frozen probe scores 95.0/95.4/
95.3/96.0\% at depths 3/6/9/12; the rule selects depth 3 (within 1pt of full).
\emph{All five datasets independently select depth 3.} Empirically fine-tuning
each depth confirms the probe: depth 3 matches or beats full depth on intent
accuracy on ATIS/SNIPS while being $\sim$4$\times$ faster.

\begin{table}[t]\centering
\caption{Multi-seed intent accuracy (mean$\pm$std, 3 seeds; MASSIVE 2 seeds).
The probe-selected depth-3 is statistically indistinguishable from full depth on
ATIS/SNIPS/CLINC150; MASSIVE and BANKING77 show a small, seed-robust gain from
depth. Every depth-3 config is $\sim$4$\times$ faster.}
\begin{tabular}{lrrc}
\toprule
Dataset & depth 3 & depth 12 & indistinguishable?\\
\midrule
ATIS & 97.95$\pm$0.14 & 97.84$\pm$0.29 & yes\\
SNIPS & 98.19$\pm$0.18 & 98.14$\pm$0.00 & yes\\
CLINC150 & 78.06$\pm$0.91 & 80.44$\pm$1.68 & yes (overlap)\\
MASSIVE & 73.49$\pm$0.32 & 74.50$\pm$0.12 & no ($\sim$1pt)\\
BANKING77 & 78.74$\pm$0.29 & 81.32$\pm$0.81 & no ($\sim$2.6pt)\\
\bottomrule
\end{tabular}
\end{table}

\begin{table}[t]\centering
\caption{Depth Pareto (ours). Accuracy peaks around depth 6--9 then plateaus;
latency is $\sim$linear in depth. Intent accuracy / slot F1 / P50 ms.}
\small
\begin{tabular}{lrrrr}
\toprule
Dataset & depth 3 & depth 6 & depth 9 & depth 12\\
\midrule
ATIS & 97.95/91.0/5.1 & 97.78/92.3/10.9 & 98.12/92.2/15.8 & 97.44/91.9/19.5\\
SNIPS & 98.00/82.2/4.9 & 98.29/84.5/9.5 & 98.43/84.9/15.7 & 98.14/84.7/18.9\\
MASSIVE & 73.81/50.6/2.7 & 74.98/55.7/5.3 & 75.62/55.1/7.6 & 74.61/54.3/10.5\\
CLINC150 & 79.02/--/2.6 & 81.44/--/5.1 & 81.93/--/7.7 & 82.00/--/10.1\\
BANKING77 & 78.34/--/2.9 & 80.91/--/5.4 & 81.79/--/9.1 & 80.78/--/10.7\\
\bottomrule
\end{tabular}
\end{table}

\textbf{5.2 Reliability (C2).} Generative parse-failure rate tracks slot-schema
complexity: 46.7\% on ATIS (101 tag types), 3.3\% on SNIPS (72), 0\% on
intent-only CLINC150; generative latency likewise scales with output length
(359/243/83\,ms). The discriminative model is flat ($\sim$5--7\,ms) and always
valid, independent of schema complexity.

\textbf{5.3 CRF ablation.} With a length-normalized CRF loss, the CRF slot head
gives a small consistent gain over softmax at depth 3 (ATIS $+0.19$, SNIPS
$+0.27$, MASSIVE $+0.09$ F1); it does not close the gap to the bidirectional
encoder, indicating the encoder's slot advantage is attention directionality,
not tag-transition modeling.

\textbf{5.4 Fair head-to-head.} At a matched 5000-example budget, ours (depth 3)
vs.\ DistilBERT: ATIS 97.95 vs 97.95 intent (tie), 91.0 vs 94.5 slot F1; the
encoder leads on slots and on fine-grained intent (CLINC150 82.0 vs 92.96). Our
recipe's value is converting an existing decoder into a fast discriminative
model, not beating a purpose-built encoder.

\section{Negative Results: CPU Inference Optimizations}
On the depth-3 ATIS model (identical weights; P50 ms): FP32 5.13; dynamic INT8
5.02 (no gain, $-0.28$ slot F1, and fails to run without an explicit quantized
engine); \texttt{torch.compile} 5.39 (\emph{slower}, overhead dominates at this
scale); ONNX Runtime 1.14 ($4.5\times$ faster, identical accuracy, but the
modern exporter fails on the decoder and the legacy exporter requires
tensor-only output surgery). Takeaway: for shallow closed-set NLU the win came
from \emph{pruning + discriminative single-pass}, not post-hoc kernel/quant
tricks; ONNX graph execution is an orthogonal, stackable win with export
friction.

\section{Discussion and Limitations}
Closed-set assumption; open-set/novel-intent detection is out of scope.
Generation is still needed for free-form tasks. Results use a single 124M
backbone and a single CPU; a second model size and rotation-INT8 with fused
kernels are future work. The fine-grained gap on BANKING77 shows depth
genuinely helps for many same-domain intents---probe-guided pruning trades a
small, quantified accuracy cost there for a $4\times$ speedup.

\section{Conclusion}
Generate less, classify more: a cheap probe selects a shallow depth, and
discriminative heads deliver fast, reliable, single-pass intent + slot
prediction from an existing generative SLM.

\section*{Reproducibility Statement}
All datasets are public. Code, exact commands, model revisions, seeds, and the
CPU specification are released via an anonymized repository. Every table is
regenerated by a single aggregation script from per-run JSON.

\section*{AI-Use Statement}
The authors used an AI coding assistant to help implement experiment scaffolding,
run the training/evaluation harness, and draft prose. All experimental design,
result interpretation, and claims were verified by the authors; all reported
numbers were produced by the released code on public benchmarks.

\bibliographystyle{plainnat}
\bibliography{references}
\end{document}
